\chapter{Resultados}

Durante os testes práticos de geração de código com inteligência artificial, constatou-se a constante necessidade de iterações manuais para se obter um código compilável ou logicamente correto. Essa característica evidenciou-se fortemente nos testes com JavaScript, que, por não possuir um compilador rigoroso e ser interpretado em tempo de execução, permitiu que a ferramenta gerasse lógicas falhas mascaradas por uma sintaxe aparentemente correta. 

Nos ensaios realizados, os LLMs raramente conseguiram gerar um código seguro e sem erros de execução na primeira tentativa. A inteligência artificial demonstrou capacidade teórica para lidar com a gestão de memória (no caso da linguagem C), casos de borda e erros de lógica; contudo, na prática, ela apenas aplicava essas correções quando o programador humano percebia a omissão e a corrigia ativamente via \textit{prompt}. 

Esse cenário comprovou a manifestação do viés de automação no ambiente de desenvolvimento: caso o programador confie excessivamente na ferramenta ou não possua conhecimentos aprofundados para guiar a IA rumo a um código seguro, os defeitos são integrados silenciosamente ao sistema. A seguir, são detalhados os resultados das avaliações qualitativas, divididos pelas linguagens e comparados sob a ótica da norma ISO/IEC 25010.

